\documentclass[12pt,a4paper]{article}

%% ── Packages ──────────────────────────────────────────────────────────────
\usepackage[T1]{fontenc}
\usepackage[utf8]{inputenc}
\usepackage{lmodern}
\usepackage{amsmath,amssymb,amsthm}
\usepackage{mathtools}
\usepackage{geometry}
\usepackage{enumitem}
\usepackage{hyperref}
\usepackage{xcolor}
\usepackage{mdframed}
\usepackage{booktabs}
\usepackage{array}
\usepackage{fancyhdr}
\usepackage{titlesec}
\usepackage{microtype}
\usepackage{bm}
\usepackage{bbm}

%% ── Page geometry ─────────────────────────────────────────────────────────
\geometry{top=2.5cm, bottom=2.5cm, left=2.8cm, right=2.8cm}
\setlength{\headheight}{15pt}

%% ── Colors ────────────────────────────────────────────────────────────────
\definecolor{metu-red}{RGB}{160,0,0}
\definecolor{proof-bg}{RGB}{248,248,252}
\definecolor{result-bg}{RGB}{240,248,240}
\definecolor{note-bg}{RGB}{255,252,235}
\definecolor{dark-gray}{RGB}{50,50,50}

%% ── Header / Footer ───────────────────────────────────────────────────────
\pagestyle{fancy}
\fancyhf{}
\fancyhead[L]{\small\color{metu-red}\textbf{EE 533} – Information Theory}
\fancyhead[R]{\small\color{dark-gray}Assignment 1 Solutions}
\fancyfoot[C]{\small\thepage}
\renewcommand{\headrulewidth}{0.4pt}
\renewcommand{\footrulewidth}{0.2pt}

%% ── Section styling ───────────────────────────────────────────────────────
\titleformat{\section}[block]
  {\large\bfseries\color{metu-red}}
  {\thesection.}{0.5em}{}[{\color{metu-red}\titlerule[0.6pt]}]

%% ── Theorem environments ──────────────────────────────────────────────────
\theoremstyle{plain}
\newtheorem{theorem}{Theorem}[section]
\newtheorem{lemma}[theorem]{Lemma}
\theoremstyle{definition}
\newtheorem{definition}[theorem]{Definition}
\newtheorem{example}[theorem]{Example}
\theoremstyle{remark}
\newtheorem*{remark}{Remark}
\newtheorem*{note}{Note}

%% ── Custom framed environments ────────────────────────────────────────────
\newmdenv[
  backgroundcolor=result-bg,
  linecolor=green!50!black,
  linewidth=1.2pt,
  roundcorner=4pt,
  innertopmargin=6pt,
  innerbottommargin=6pt,
  innerleftmargin=10pt,
  innerrightmargin=10pt
]{resultbox}

\newmdenv[
  backgroundcolor=note-bg,
  linecolor=orange!70!black,
  linewidth=1pt,
  roundcorner=4pt,
  innertopmargin=6pt,
  innerbottommargin=6pt,
  innerleftmargin=10pt,
  innerrightmargin=10pt
]{notebox}

\newmdenv[
  backgroundcolor=proof-bg,
  linecolor=gray!60,
  linewidth=0.8pt,
  roundcorner=3pt,
  innertopmargin=6pt,
  innerbottommargin=6pt,
  innerleftmargin=10pt,
  innerrightmargin=10pt
]{proofbox}

%% ── Math macros ───────────────────────────────────────────────────────────
\newcommand{\HH}{\mathrm{H}}
\newcommand{\hb}{\mathrm{h}}          % binary entropy
\newcommand{\II}{\mathrm{I}}
\newcommand{\EE}{\mathbf{E}}
\newcommand{\PP}{\mathbf{P}}
\newcommand{\supp}{\mathrm{supp}}
\newcommand{\expeq}{\;\dot{=}\;}       % first-order exponent equality

%% ── Hyperref setup ────────────────────────────────────────────────────────
\hypersetup{
  colorlinks=true,
  linkcolor=metu-red,
  urlcolor=metu-red,
  citecolor=metu-red
}

%% ═══════════════════════════════════════════════════════════════════════════
\begin{document}
%% ═══════════════════════════════════════════════════════════════════════════

%% ── Title Page ────────────────────────────────────────────────────────────
\begin{titlepage}
  \centering
  \vspace*{3cm}
  {\Large\color{metu-red}\textbf{Middle East Technical University}\\[0.3em]
   Department of Electrical \& Electronics Engineering}\\[2em]
  \rule{\linewidth}{1pt}\\[1em]
  {\Huge\bfseries EE 533 – Information Theory}\\[0.8em]
  {\Large\bfseries Assignment 1 – Solutions}\\[1em]
  \rule{\linewidth}{1pt}\\[3em]
  {\large
    \begin{tabular}{ll}
      \textbf{Instructor:} & Barış Nakiboğlu \\[0.5em]
      \textbf{Semester:}   & Spring 2025 \\
    \end{tabular}
  }\\[4em]
  {\normalsize
    \begin{tabular}{>{\bfseries}l l}
      Problems: & 2.2,\; 2.3,\; 2.10,\; 2.39,\; 3.3 \\
      Textbook: & Cover \& Thomas, \textit{Elements of Information Theory},\\
                & 2\textsuperscript{nd} ed.\ (Wiley, 2006)
    \end{tabular}
  }
  \vfill
  {\small\color{dark-gray} All logarithms are base~2 (bits) unless otherwise noted.}
\end{titlepage}

%% ── Table of Contents ─────────────────────────────────────────────────────
\tableofcontents
\newpage

%%============================================================================
\section{Problem 2.2 – Entropy of Functions}
%%============================================================================

\subsection*{Problem Statement}

Let $X$ be a random variable taking on a \emph{finite} number of values.
What is the (general) inequality relation between $\HH(X)$ and $\HH(Y)$ if:

\begin{enumerate}[label=\textbf{(\alph*)}]
  \item $Y = 2^X$
  \item $Y = \cos X$
\end{enumerate}

\subsection*{Background}

We recall two fundamental results from lecture (Week~1).

\begin{theorem}[Invariance under relabelling]\label{thm:inv}
  For any one-to-one (injective) function $f$,
  \[
    \HH(f(X)) = \HH(X).
  \]
\end{theorem}

\begin{theorem}[Data-processing / Function inequality]\label{thm:func}
  For any (possibly non-injective) function $g$,
  \[
    \HH(g(X)) \;\le\; \HH(X),
  \]
  with equality if and only if $g$ is injective on $\supp(p_X)
  := \{x : p_X(x)>0\}$.
\end{theorem}

\subsection*{Solution}

%% ── Part (a) ──────────────────────────────────────────────────────────────
\subsubsection*{(a) $Y = 2^X$}

\paragraph{Step 1 – Show $f(x) = 2^x$ is injective.}

Suppose $2^{x_1} = 2^{x_2}$ for some $x_1, x_2$ in the support of $p_X$.
Taking $\log_2$ of both sides (which is the inverse of $2^x$, hence well-defined since $2^x>0$ for all $x$):
\[
  2^{x_1} = 2^{x_2}
  \;\Longrightarrow\;
  \log_2\!\left(2^{x_1}\right) = \log_2\!\left(2^{x_2}\right)
  \;\Longrightarrow\;
  x_1 = x_2.
\]
Hence $f(x)=2^x$ is \textbf{injective} on any set of real numbers.

\paragraph{Step 2 – Apply Theorem~\ref{thm:inv}.}

Because $f$ is injective, Theorem~\ref{thm:inv} gives immediately:

\begin{resultbox}
\[
  \boxed{\HH(Y) = \HH(X).}
\]
\end{resultbox}

\begin{note}
Knowing $Y = 2^X$ is \emph{exactly} as informative as knowing $X$: given $Y$
one can recover $X = \log_2 Y$ uniquely. No information is created or destroyed
by a bijective relabelling.
\end{note}

%% ── Part (b) ──────────────────────────────────────────────────────────────
\subsubsection*{(b) $Y = \cos X$}

\paragraph{Step 1 – Cosine is not injective in general.}

The cosine function satisfies $\cos(\theta) = \cos(-\theta)$ for all $\theta$.
Therefore if $\supp(p_X)$ contains both $\theta$ and $-\theta$ for some
$\theta\neq 0$, then $g(x)=\cos(x)$ is not injective on $\supp(p_X)$, and
\emph{information is lost} in the mapping $X\mapsto Y$.

\paragraph{Step 2 – Formal chain-rule argument.}

Since $Y = g(X)$ is a deterministic function of $X$, we have $H(Y\mid X)=0$,
and hence $H(X,Y)=H(X)$.  By the chain rule:
\[
  \HH(X) \;=\; \HH(X,Y) \;=\; \HH(Y) + \HH(X \mid Y).
\]
Since $\HH(X\mid Y) \ge 0$, it follows that

\begin{resultbox}
\[
  \boxed{\HH(Y) \;\le\; \HH(X).}
\]
\end{resultbox}

\paragraph{Step 3 – Equality condition.}

The proof shows:
\[
  \HH(Y) = \HH(X)
  \;\Longleftrightarrow\;
  \HH(X\mid Y) = 0
  \;\Longleftrightarrow\;
  X \text{ is a (deterministic) function of } Y
  \;\Longleftrightarrow\;
  g\text{ is injective on }\supp(p_X).
\]

\begin{notebox}
\textbf{Example where equality holds.}
If $\supp(p_X) \subseteq [0,\pi]$, then cosine is strictly decreasing (hence
injective) on that interval, so $\HH(\cos X)=\HH(X)$.  The inequality is
therefore strict only when $g$ causes genuine ``collisions'' in the support.
\end{notebox}

\begin{notebox}
\textbf{Notation.}
Throughout this problem we use $f(x)=2^x$ for part~(a) and $g(x)=\cos x$ for
part~(b) to keep the two functions notationally distinct.
\end{notebox}

\newpage
%%============================================================================
\section{Problem 2.3 – Minimum Entropy}
%%============================================================================

\subsection*{Problem Statement}

What is the minimum value of
\[
  \HH(p_1,\dots,p_n) \;=\; \HH(\mathbf{p})
\]
as $\mathbf{p}$ ranges over the set of all $n$-dimensional probability vectors?
Find all $\mathbf{p}$'s that achieve this minimum.

\subsection*{Background}

\begin{definition}[Shannon Entropy]
  For a discrete random variable $X$ with pmf $p_X$,
  \[
    \HH(X) \;:=\; -\sum_{x} p_X(x)\,\log p_X(x) \;=\; \EE\!\left[\log\frac{1}{p_X(X)}\right],
  \]
  where we adopt the convention $0\log 0 := 0$.
  \emph{Motivation of the convention:} By L'Hôpital's rule,
  $\lim_{p\to 0^+} p\log p = 0$,
  so the convention preserves continuity of each summand at $p_i=0$.
\end{definition}

\subsection*{Solution}

\begin{theorem}
  \[
    \min_{\mathbf{p}\in\Delta_n}\; \HH(\mathbf{p}) \;=\; 0,
  \]
  where $\Delta_n := \{\mathbf{p}\in[0,1]^n : \sum_{i=1}^n p_i = 1\}$.
  The minimum is achieved \emph{if and only if} $\mathbf{p} = \mathbf{e}_k$
  for some $k\in\{1,\dots,n\}$ (a point-mass / deterministic distribution).
\end{theorem}

\begin{proof}
We proceed in three steps.

\textbf{Step 1 – Lower bound $\HH(\mathbf{p})\ge 0$.}

For every $p_i \in [0,1]$ we have $\log p_i \le 0$, so
$-p_i\log p_i \ge 0$ (with the convention $0\cdot\log 0 = 0$).
Since the entropy is a sum of non-negative terms,
\[
  \HH(\mathbf{p}) = -\sum_{i=1}^n p_i\log p_i \;\ge\; 0.
\]

\textbf{Step 2 – Achievability.}

Set $p_k = 1$ for some fixed $k$ and $p_i = 0$ for all $i\neq k$.  Then
\[
  \HH(\mathbf{p}) = -1\cdot\log 1 - \sum_{i\neq k} 0\cdot\log 0 = 0 - 0 = 0.
\]

\textbf{Step 3 – Uniqueness.}

Suppose $\HH(\mathbf{p}) = 0$.  Then every summand $-p_i\log p_i = 0$.
Since $-p\log p = 0$ if and only if $p\in\{0,1\}$ (the function $-p\log p$ is
strictly positive for $p\in(0,1)$), we conclude $p_i\in\{0,1\}$ for all $i$.
Combined with $\sum_{i=1}^n p_i = 1$, exactly one $p_k = 1$ and all others
equal $0$.  There are exactly $n$ such vectors: $\mathbf{e}_1,\dots,\mathbf{e}_n$.
\end{proof}

\begin{resultbox}
\[
  \boxed{\min_{\mathbf{p}\in\Delta_n} \HH(\mathbf{p}) \;=\; 0,}
\]
achieved \textbf{if and only if} $\mathbf{p}\in\{\mathbf{e}_1,\dots,\mathbf{e}_n\}$,
i.e., the distribution is a point mass (the outcome is \emph{deterministic}).
\end{resultbox}

\begin{note}
The maximum is $\HH(\mathbf{p})\le \log n$, attained uniquely at the uniform
distribution $p_i = 1/n$ for all $i$.  This follows from log-sum inequality /
Jensen's inequality applied to the strictly concave function $-p\log p$.
\end{note}

\newpage
%%============================================================================
\section{Problem 2.10 – Entropy of a Disjoint Mixture}
%%============================================================================

\subsection*{Problem Statement}

Let $X_1$ and $X_2$ be discrete random variables drawn according to pmfs
$p_1(\cdot)$ and $p_2(\cdot)$ over \emph{disjoint} alphabets
$\mathcal{X}_1 = \{1,\dots,m\}$ and $\mathcal{X}_2 = \{m+1,\dots,n\}$.
Define the mixture variable
\[
  X = \begin{cases}
    X_1 & \text{with probability } \alpha,\\
    X_2 & \text{with probability } 1-\alpha,
  \end{cases}
  \qquad \alpha\in(0,1).
\]

\begin{enumerate}[label=\textbf{(\alph*)}]
  \item Find $\HH(X)$ in terms of $\HH(X_1)$, $\HH(X_2)$, and $\alpha$.
  \item Maximize over $\alpha$ to show $2^{\HH(X)} \le 2^{\HH(X_1)} + 2^{\HH(X_2)}$
        and interpret using the notion that $2^{\HH(X)}$ is the effective alphabet size.
\end{enumerate}

\subsection*{Background}

Denote the \emph{binary entropy function}
$\hb(\alpha) := -\alpha\log\alpha - (1-\alpha)\log(1-\alpha)$ for
$\alpha\in(0,1)$.

\subsection*{Solution}

%% ── Part (a) ──────────────────────────────────────────────────────────────
\subsubsection*{(a) Expression for $\HH(X)$}

\paragraph{Step 1 – PMF of $X$.}

Because $\mathcal{X}_1\cap\mathcal{X}_2=\emptyset$, the pmf of $X$ is:
\[
  p_X(k) = \begin{cases}
    \alpha\, p_1(k)       & k \in \mathcal{X}_1,\\
    (1-\alpha)\,p_2(k)   & k \in \mathcal{X}_2.
  \end{cases}
\]

\paragraph{Step 2 – Split the entropy sum.}

\begin{align*}
  \HH(X)
    &= -\sum_{k\in\mathcal{X}_1}\!\alpha p_1(k)\log\!\bigl(\alpha p_1(k)\bigr)
       -\sum_{k\in\mathcal{X}_2}\!(1-\alpha)p_2(k)\log\!\bigl((1-\alpha)p_2(k)\bigr).
\end{align*}

\paragraph{Step 3 – Expand $\log(\alpha p_i) = \log\alpha + \log p_i$.}

\begin{align*}
  \HH(X)
    &= -\sum_{k\in\mathcal{X}_1}\alpha p_1(k)\bigl[\log\alpha + \log p_1(k)\bigr]
       -\sum_{k\in\mathcal{X}_2}(1-\alpha)p_2(k)\bigl[\log(1-\alpha)+\log p_2(k)\bigr]\\[4pt]
    &= \underbrace{-\alpha\log\alpha}_{}\!\underbrace{\sum_{k}p_1(k)}_{=\,1}
       +\;\alpha\!\underbrace{\Bigl(-\sum_{k}p_1(k)\log p_1(k)\Bigr)}_{=\;\HH(X_1)}\\[2pt]
    &\quad -\underbrace{(1-\alpha)\log(1-\alpha)}_{}\!\underbrace{\sum_{k}p_2(k)}_{=\,1}
       +\;(1-\alpha)\!\underbrace{\Bigl(-\sum_{k}p_2(k)\log p_2(k)\Bigr)}_{=\;\HH(X_2)}.
\end{align*}

\paragraph{Step 4 – Collect terms.}

\begin{resultbox}
\[
  \boxed{\HH(X) \;=\; \hb(\alpha) + \alpha\,\HH(X_1) + (1-\alpha)\,\HH(X_2).}
\]
\end{resultbox}

\begin{notebox}
\textbf{Interpretation.}
The entropy of the mixture decomposes into two parts:
\begin{enumerate}
  \item $\hb(\alpha)$: the uncertainty about \emph{which component} was selected.
  \item $\alpha\HH(X_1)+(1-\alpha)\HH(X_2)$: the expected uncertainty
        \emph{within} the selected component.
\end{enumerate}
This is the chain rule in disguise: $\HH(X) = \HH(B) + \HH(X\mid B)$,
where $B$ is the Bernoulli$(\alpha)$ selector.
\end{notebox}

%% ── Part (b) ──────────────────────────────────────────────────────────────
\subsubsection*{(b) Upper bound $2^{\HH(X)} \le 2^{\HH(X_1)} + 2^{\HH(X_2)}$}

\paragraph{Step 1 – Rewrite in terms of $a$ and $b$.}

Let $a := 2^{\HH(X_1)}$ and $b := 2^{\HH(X_2)}$.  Using the result of part~(a):
\begin{align*}
  2^{\HH(X)}
    &= 2^{\;\hb(\alpha) \;+\; \alpha\,\HH(X_1) \;+\; (1-\alpha)\,\HH(X_2)}\\[4pt]
    &= 2^{\hb(\alpha)}\cdot 2^{\,\alpha\,\HH(X_1)}\cdot 2^{\,(1-\alpha)\,\HH(X_2)}\\[4pt]
    &= 2^{\hb(\alpha)}\cdot \bigl(2^{\HH(X_1)}\bigr)^{\!\alpha}
                           \cdot\bigl(2^{\HH(X_2)}\bigr)^{\!1-\alpha}\\[4pt]
    &= 2^{\hb(\alpha)}\cdot a^{\alpha}\cdot b^{1-\alpha}.
\end{align*}

\paragraph{Step 2 – Maximize $\HH(X)$ over $\alpha$.}

Differentiating with respect to $\alpha$ and setting equal to zero:
\begin{align*}
  \frac{d}{d\alpha}\HH(X)
    &= \frac{d}{d\alpha}\Bigl[\hb(\alpha) + \alpha\log a + (1-\alpha)\log b\Bigr]\\[4pt]
    &= \frac{d\hb}{d\alpha} + \log a - \log b\\[4pt]
    &= \log\!\frac{1-\alpha}{\alpha} + \log a - \log b \;=\; 0,
\end{align*}
where we used $\hb'(\alpha) = \log\!\frac{1-\alpha}{\alpha}$.  Solving:
\[
  \log\!\frac{1-\alpha}{\alpha} = \log\!\frac{b}{a}
  \;\Longrightarrow\;
  \frac{1-\alpha}{\alpha} = \frac{b}{a}
  \;\Longrightarrow\;
  \boxed{\alpha^* = \frac{a}{a+b} = \frac{2^{\HH(X_1)}}{2^{\HH(X_1)}+2^{\HH(X_2)}}.}
\]
Since $\HH(X)$ is strictly concave in $\alpha$ (as a sum of concave functions),
this critical point is the \emph{global maximum}.

\paragraph{Step 3 – Evaluate $2^{\HH(X)}$ at $\alpha^*$.}

At $\alpha^*$:
\[
  \hb(\alpha^*)
    = -\frac{a}{a+b}\log\frac{a}{a+b} - \frac{b}{a+b}\log\frac{b}{a+b}
    = \frac{a}{a+b}\log\frac{a+b}{a} + \frac{b}{a+b}\log\frac{a+b}{b}.
\]
Therefore:
\begin{align*}
  2^{\HH(X^*)}
    &= 2^{\hb(\alpha^*)}\cdot a^{\tfrac{a}{a+b}}\cdot b^{\tfrac{b}{a+b}}\\[6pt]
    &= \left(\frac{a+b}{a}\right)^{\!\frac{a}{a+b}}
       \!\cdot\left(\frac{a+b}{b}\right)^{\!\frac{b}{a+b}}
       \!\cdot\; a^{\tfrac{a}{a+b}}\cdot b^{\tfrac{b}{a+b}}\\[6pt]
    &= (a+b)^{\tfrac{a}{a+b}}\cdot(a+b)^{\tfrac{b}{a+b}}
     \;=\; (a+b)^{\tfrac{a+b}{a+b}}
     \;=\; a+b.
\end{align*}

\paragraph{Step 4 – Conclude the bound.}

Since $\alpha^*$ is the \emph{maximizer} of $\HH(X)$, for \emph{every}
$\alpha\in(0,1)$:
\[
  \HH(X) \;\le\; \HH(X^*),
  \quad\text{hence}\quad
  2^{\HH(X)} \;\le\; 2^{\HH(X^*)}.
\]

\begin{resultbox}
\[
  \boxed{2^{\HH(X)} \;\le\; 2^{\HH(X_1)} + 2^{\HH(X_2)},}
\]
with equality if and only if $\alpha = \alpha^* = \dfrac{2^{\HH(X_1)}}{2^{\HH(X_1)}+2^{\HH(X_2)}}$.
\end{resultbox}

\paragraph{Interpretation.}
The quantity $2^{\HH(X)}$ represents the \emph{effective alphabet size}: the
number of roughly equally likely outcomes.  The bound states that the effective
alphabet size of a mixture is at most the \emph{sum} of the effective alphabet
sizes of the components.  The optimal mixing weight $\alpha^*$ is proportional
to the effective size of each component, balancing the two distributions to
maximize total uncertainty.

\newpage
%%============================================================================
\section{Problem 2.39 – Entropy and Pairwise Independence}
%%============================================================================

\subsection*{Problem Statement}

Let $X$, $Y$, $Z$ be three binary Bernoulli$(\tfrac{1}{2})$ random variables
that are \emph{pairwise independent}:
\[
  \II(X;Y) = \II(X;Z) = \II(Y;Z) = 0.
\]

\begin{enumerate}[label=\textbf{(\alph*)}]
  \item Under this constraint, what is the \emph{minimum} value of $\HH(X,Y,Z)$?
  \item Give an example achieving this minimum.
\end{enumerate}

\subsection*{Background}

\begin{definition}[Mutual Information]
  $\II(X;Y) := \HH(X) + \HH(Y) - \HH(X,Y)$.
  Hence $\II(X;Y) = 0 \Longleftrightarrow \HH(X,Y) = \HH(X) + \HH(Y)$
  (i.e., $X$ and $Y$ are \emph{independent}).
\end{definition}

\subsection*{Solution}

%% ── Part (a) ──────────────────────────────────────────────────────────────
\subsubsection*{(a) Minimum value of $\HH(X,Y,Z)$}

\paragraph{Step 1 – Marginal entropies.}

Since each variable is Bernoulli$(\tfrac{1}{2})$:
\[
  \HH(X) = \HH(Y) = \HH(Z) = \hb\!\left(\tfrac{1}{2}\right) = 1\text{ bit}.
\]

\paragraph{Step 2 – Joint entropy of each pair.}

Pairwise independence $\II(X;Y)=0$ means, by definition:
\[
  \HH(X,Y) \;=\; \HH(X) + \HH(Y) - \underbrace{\II(X;Y)}_{=\,0}
  \;=\; \HH(X) + \HH(Y) \;=\; 2\text{ bits.}
\]
Likewise $\HH(X,Z) = \HH(Y,Z) = 2$ bits.

\paragraph{Step 3 – Lower bound on $\HH(X,Y,Z)$.}

By the \emph{chain rule for entropy}:
\[
  \HH(X,Y,Z)
    \;=\; \HH(X,Y) + \HH(Z\mid X,Y)
    \;=\; 2 + \HH(Z\mid X,Y).
\]
Since $\HH(Z\mid X,Y) \ge 0$ always,
\[
  \HH(X,Y,Z) \;\ge\; 2\text{ bits.}
\]

\paragraph{Step 4 – Characterisation of equality.}

The lower bound is achieved ($\HH(X,Y,Z)=2$) \emph{if and only if}
\[
  \HH(Z\mid X,Y) = 0 \;\Longleftrightarrow\; Z \text{ is a deterministic function of } (X,Y).
\]
We must verify there \emph{exists} such a function satisfying all pairwise
independence constraints — this is done in part~(b).

\begin{resultbox}
\[
  \boxed{\min \HH(X,Y,Z) = 2\text{ bits.}}
\]
\end{resultbox}

%% ── Part (b) ──────────────────────────────────────────────────────────────
\subsubsection*{(b) Achieving example: $Z = X \oplus Y$}

\paragraph{Construction.}

Let $X$ and $Y$ be i.i.d.\ Bernoulli$(\tfrac{1}{2})$, and define
$Z := X \oplus Y$ (XOR, equivalently addition mod~2).

\paragraph{Verification 1 – $Z$ is Bernoulli$(\tfrac{1}{2})$.}
\[
  \PP(Z=0) = \PP(X=Y) = \PP(X=0,Y=0)+\PP(X=1,Y=1) = \tfrac{1}{4}+\tfrac{1}{4} = \tfrac{1}{2}.\quad\checkmark
\]

\paragraph{Verification 2 – Pairwise independence.}

Since $X,Y$ are i.i.d., $\II(X;Y)=0$ by assumption.

For $X\perp Z$: for any $(x,z)\in\{0,1\}^2$, the unique value $y=z\oplus x$
satisfies $(X=x)\cap(Z=z) = (X=x)\cap(Y=z\oplus x)$, so
\[
  \PP(X=x,\,Z=z)
    = \PP(X=x,\,Y=z\oplus x)
    = \underbrace{\PP(X=x)}_{1/2}\cdot\underbrace{\PP(Y=z\oplus x)}_{1/2}
    = \frac{1}{4}
    = \PP(X=x)\cdot\PP(Z=z). \quad\checkmark
\]
For $Y\perp Z$: since $Z = X\oplus Y = Y\oplus X$ (XOR is commutative),
for any $(y,z)\in\{0,1\}^2$ the unique value $x = z\oplus y$ gives
$(Y=y)\cap(Z=z) = (X=z\oplus y)\cap(Y=y)$, so
\[
  \PP(Y=y,\,Z=z)
    = \PP(X=z\oplus y,\,Y=y)
    = \underbrace{\PP(X=z\oplus y)}_{1/2}\cdot\underbrace{\PP(Y=y)}_{1/2}
    = \frac{1}{4}
    = \PP(Y=y)\cdot\PP(Z=z). \quad\checkmark
\]

\paragraph{Verification 3 – Entropy.}

$Z$ is fully determined by $(X,Y)$, so $\HH(Z\mid X,Y)=0$, and:
\[
  \HH(X,Y,Z) = \HH(X,Y) + \HH(Z\mid X,Y) = 2 + 0 = 2\text{ bits.} \quad\checkmark
\]
The lower bound of Step~4 is therefore \emph{tight}.

\begin{notebox}
\textbf{Non-mutual independence.}
Although $X$, $Y$, $Z$ are pairwise independent, they are \textbf{not} mutually
independent: $\HH(X,Y,Z) = 2 < 3 = \HH(X)+\HH(Y)+\HH(Z)$.  Equivalently,
knowing both $X$ and $Y$ completely determines $Z$ — a clear dependence.
This is a classical counterexample showing that
\emph{pairwise independence does not imply mutual independence}.
\end{notebox}

\newpage
%%============================================================================
\section{Problem 3.3 – Piece of Cake}
%%============================================================================

\subsection*{Problem Statement}

A cake is sliced roughly in half; the largest piece is always chosen and the
other discarded.  Each random cut creates pieces in proportions
\[
  P = \begin{cases}
    \bigl(\tfrac{2}{3},\,\tfrac{1}{3}\bigr) & \text{with probability } \tfrac{3}{4},\\[4pt]
    \bigl(\tfrac{3}{5},\,\tfrac{2}{5}\bigr) & \text{with probability } \tfrac{1}{4}.
  \end{cases}
\]
How large, \emph{to first order in the exponent}, is the piece of cake after
$n$ cuts?

\subsection*{Background}

\begin{theorem}[Weak Law of Large Numbers]
  Let $W_1,W_2,\dots$ be i.i.d.\ with $\EE[|W_1|]<\infty$.  Then
  $\frac{1}{n}\sum_{k=1}^n W_k \xrightarrow{\;\PP\;} \EE[W_1]$ as $n\to\infty$.
\end{theorem}

\begin{definition}[First-order exponent equality]
  We write $S_n \\expeq 2^{nR}$ to mean $\frac{1}{n}\log_2 S_n \to R$ in
  probability as $n\to\infty$.
\end{definition}

\subsection*{Solution}

\paragraph{Step 1 – Model.}

At each cut $i$, the surviving (largest) piece is multiplied by
\[
  M_i = \begin{cases}
    \tfrac{2}{3} & \text{with probability } \tfrac{3}{4},\\[2pt]
    \tfrac{3}{5} & \text{with probability } \tfrac{1}{4},
  \end{cases}
\]
with $M_1, M_2, \dots$ i.i.d.  The cake size after $n$ cuts (starting from
size~$1$) is
\[
  S_n = \prod_{i=1}^{n} M_i.
\]

\paragraph{Step 2 – Take logarithms.}

\[
  \frac{1}{n}\log_2 S_n
    = \frac{1}{n}\sum_{i=1}^{n}\log_2 M_i.
\]

\paragraph{Step 3 – Verify the LLN condition.}

Since $M_i\in\{\tfrac{2}{3},\tfrac{3}{5}\}$ are bounded and positive,
$|\log_2 M_i|$ is bounded, so $\EE[|\log_2 M_1|]<\infty$ trivially.
By the WLLN:
\[
  \frac{1}{n}\log_2 S_n \;\xrightarrow{\;\PP\;}\; \EE[\log_2 M]
  \quad\text{as }n\to\infty.
\]

\paragraph{Step 4 – Compute $\EE[\log_2 M]$.}

\begin{align*}
  \EE[\log_2 M]
    &= \frac{3}{4}\log_2\!\frac{2}{3} + \frac{1}{4}\log_2\!\frac{3}{5}\\[4pt]
    &= \frac{3}{4}\bigl(\log_2 2 - \log_2 3\bigr)
       +\frac{1}{4}\bigl(\log_2 3 - \log_2 5\bigr)\\[4pt]
    &= \frac{3}{4} - \frac{3}{4}\log_2 3
       + \frac{1}{4}\log_2 3 - \frac{1}{4}\log_2 5\\[4pt]
    &= \frac{3}{4} - \frac{1}{2}\log_2 3 - \frac{1}{4}\log_2 5.
\end{align*}

\noindent\textit{Numerical evaluation:}
$\log_2 3 \approx 1.58496$,\; $\log_2 5 \approx 2.32193$.
\[
  \EE[\log_2 M]
    \approx 0.750 - \tfrac{1}{2}(1.585) - \tfrac{1}{4}(2.322)
    = 0.750 - 0.793 - 0.581
    = -0.624.
\]

\paragraph{Step 5 – Size to first order in the exponent.}

By the definition of $\\expeq$:
\[
  S_n
    \;\\expeq\;
    2^{\,n\,\EE[\log_2 M]}
    \;=\;
    \left(\frac{2}{3}\right)^{\!\frac{3n}{4}}
    \cdot
    \left(\frac{3}{5}\right)^{\!\frac{n}{4}}.
\]

\paragraph{Step 6 – Closed-form exponent.}

\[
  \left(\frac{2}{3}\right)^{\!\frac{3n}{4}}
  \cdot
  \left(\frac{3}{5}\right)^{\!\frac{n}{4}}
  =
  \left[\left(\frac{2}{3}\right)^{\!3}\cdot\frac{3}{5}\right]^{\!n/4}
  =
  \left[\frac{8}{27}\cdot\frac{3}{5}\right]^{\!n/4}
  =
  \left(\frac{8}{45}\right)^{\!n/4}.
\]

\begin{notebox}
\textbf{Clarification on the exponent notation.}
The exponent $n/4$ applies to the \emph{entire} fraction $(8/45)$, not to~45 alone:
\[
  \left(\frac{8}{45}\right)^{\!n/4}
    = \frac{8^{\,n/4}}{45^{\,n/4}}
    \;\neq\;
    \frac{8}{45^{\,n/4}}.
\]
Numerical check at $n=8$:\quad
$\left(\tfrac{8}{45}\right)^{\!2} = \tfrac{64}{2025}\approx 0.0316$,\quad
whereas $8/45^{2} = 8/2025\approx 0.00395$ — a factor of~8 difference.
\end{notebox}

\begin{resultbox}
\[
  \boxed{S_n \;\\expeq\;
    \left(\frac{2}{3}\right)^{\!\frac{3n}{4}}
    \!\cdot\!
    \left(\frac{3}{5}\right)^{\!\frac{n}{4}}
    =
    \left(\frac{8}{45}\right)^{\!n/4}
    \;\approx\; 2^{-0.624\,n}.}
\]
\end{resultbox}

\begin{notebox}
\textbf{Interpretation.}
The factor $2^{-0.624}\approx 0.649$ is the \emph{geometric mean} of the cut
factors:
\[
  2^{\EE[\log_2 M]}
    = \left(\frac{2}{3}\right)^{\!3/4}\!\cdot\!\left(\frac{3}{5}\right)^{\!1/4}
    \approx 0.649.
\]
This is \emph{not} the value of any single cut (individual cuts give $2/3\approx0.667$
or $3/5=0.600$), but rather the \emph{log-average} reduction rate per cut.
In other words, $S_n\approx(0.649)^n$ in the exponential sense: each of the
$n$ cuts contributes an average of $\approx -0.624$ bits to $\log_2 S_n$.
The cake shrinks to zero exponentially fast.
\end{notebox}

\begin{notebox}
\textbf{Almost-sure convergence.}
Since $M_i$ are bounded, $\EE[(\log_2 M_1)^2]<\infty$ as well, so in fact
the Strong Law of Large Numbers gives
$\frac{1}{n}\log_2 S_n \xrightarrow{\;\text{a.s.}\;} \EE[\log_2 M]$,
a stronger statement than convergence in probability.
\end{notebox}

\end{document}
